% Iteration 5 change manifest as two structured cards: chg-7 introduces the
% prompt + tool-descriptor publish-state rules, chg-8 installs the stateful
% shell guard. Both ship together at the iteration-4 boundary as the
% iteration-5 harness. Palette and layout follow case-trajectory.tex.
\begin{figure}[t]
    \centering
    \scriptsize
    \begin{tcbraster}[raster columns=2,
                      raster equal height,
                      raster column skip=4pt,
                      raster row skip=0pt,
                      raster left skip=0pt,
                      raster right skip=0pt,
                      boxrule=0.6pt,
                      left=4pt,right=4pt,top=3pt,bottom=3pt,
                      coltitle=white]
        \begin{tcolorbox}[colback=aheSky!18,colframe=aheNavy,colbacktitle=aheNavy,
                          title={\scriptsize\textbf{chg-7, iteration 5, commit \texttt{3ba3a90}, level: prompt + tool descriptor}}]
            \textbf{\textcolor{aheNavy}{Files}} \\
            \rule{\linewidth}{0.3pt}\\[1pt]
            $\bullet$ \texttt{workspace/systemprompt.md} \\
            $\bullet$ \texttt{tool\_descriptions/run\_shell\_command.tool.yaml} \\[3pt]
            \rule{\linewidth}{0.3pt}\\
            \textbf{\textcolor{aheNavy}{What changed}} \\
            \rule{\linewidth}{0.3pt}\\[1pt]
            Appended three rules to the harness's working memory. Publish-state rule: once an evaluator-style final check passes, the resulting filesystem and service state is the deliverable surface and must not be reset to ``look clean''. Scratch-directory rule: place exploratory artifacts under \texttt{/tmp} or a scratch path the verifier ignores. Literal-output rule: for DSL, config, or script outputs with byte-level contracts, validate equality at the byte level. \\[3pt]
            \rule{\linewidth}{0.3pt}\\
            \textbf{\textcolor{aheNavy}{Failure pattern fixed}} \\
            \rule{\linewidth}{0.3pt}\\[1pt]
            Agent reached an evaluator-passing state, then issued sweeping cleanup or rewrote outputs to ``tidy up'', leaving the verifier with no deliverable. \\[3pt]
            \rule{\linewidth}{0.3pt}\\
            \textbf{\textcolor{aheNavy}{Predicted fixes}} \\
            \rule{\linewidth}{0.3pt}\\[1pt]
            4 tasks. Examples: \texttt{path-tracing}, \texttt{configure-git-webserver}, \texttt{polyglot-rust-c}, \texttt{large-scale-text-editing}.
        \end{tcolorbox}
        \begin{tcolorbox}[colback=aheTeal!12,colframe=aheBlue,colbacktitle=aheBlue,
                          title={\scriptsize\textbf{chg-8, iteration 5, commit \texttt{4e0aab9}, level: tool implementation}}]
            \textbf{\textcolor{aheBlue}{Files}} \\
            \rule{\linewidth}{0.3pt}\\[1pt]
            $\bullet$ \texttt{tools/shell\_tools/run\_shell\_command.py} \\[3pt]
            \rule{\linewidth}{0.3pt}\\
            \textbf{\textcolor{aheBlue}{What changed}} \\
            \rule{\linewidth}{0.3pt}\\[1pt]
            Installed a stateful publish-state guard inside the shell tool. After a successful evaluator-style final check, the guard parses the acceptance command for explicit file paths and roots and records them as protected. Later destructive commands that would delete a protected output or reset a protected root are intercepted before execution and returned as a targeted error. An explicit \texttt{ALLOW\_POST\_SUCCESS\_RESET} token can downgrade the block to a warning, after which the agent must re-validate before submit. \\[3pt]
            \rule{\linewidth}{0.3pt}\\
            \textbf{\textcolor{aheBlue}{Failure pattern fixed}} \\
            \rule{\linewidth}{0.3pt}\\[1pt]
            Even with the prompt rule in place, the agent still issued destructive cleanup commands after publish-state. Execution-time enforcement at the shell tool is the most direct interlock. \\[3pt]
            \rule{\linewidth}{0.3pt}\\
            \textbf{\textcolor{aheBlue}{Predicted fixes}} \\
            \rule{\linewidth}{0.3pt}\\[1pt]
            Same 4 tasks; load-bearing on \texttt{path-tracing}, whose F4 is the \texttt{rm -rf} of \texttt{/app/reconstructed.ppm}.
        \end{tcolorbox}
    \end{tcbraster}
    \caption{The two change-manifest entries written together at the iteration-4 boundary and shipped as the iteration-5 harness. \texttt{chg-7} names the publish-state rule in the system prompt and tool descriptor; \texttt{chg-8} installs the execution-time interlock inside the shell tool. The pair flips \texttt{path-tracing} on the next round.}
    \label{fig:manifest-publish-state}
\end{figure}
